\clearpage
\begin{examplebox}
\textbf{例1（最典型：已知两个交点，直接求面积）。}

设双曲线 $\dfrac{x^2}{9}-\dfrac{y^2}{4}=1$。
直线 $\ell$ 与 $y=\tfrac23 x$、$y=-\tfrac23x$
分别交于 $A(6,4)$、$B(3,-2)$，求 $S_{\triangle AOB}$。

\textbf{解：}$a=3,b=2$，$x_1=6,x_2=3$，故
\[
\boxed{S=\frac23|6\cdot3|=12}.
\]
\textbf{复核：} $2S=|6(-2)-3(4)|=24$。

\medskip
\textbf{例2（横坐标异号：绝对值不能少）。}

仍取 $a=3,b=2$，令 $A(6,4)$、$B(-3,2)$。
二点分别在 $y=\tfrac23x$、$y=-\tfrac23x$ 上，且
\[
\boxed{S=\frac23|6(-3)|=12}.
\]
若误把绝对值去掉就得到 $-12$，这只是有向面积，不是面积。

\medskip
\textbf{例3（水平截线：两截点在同侧 $y>0$）。}

取 $\ell:y=4$、$a=3,b=2$，则
$A(6,4)$、$B(-6,4)$，面积为
\[
S=\frac12\cdot AB\cdot4=\frac12\cdot12\cdot4=24;
\qquad\boxed{\frac23|6(-6)|=24}.
\]
\textbf{提醒：}水平截线根本不与 $x$ 轴相交，
所以不能默认用“$x$ 轴截距”作底。

\medskip
\textbf{例4（竖直截线：两个横坐标相等）。}

取 $\ell:x=3$、$a=3,b=2$，有 $A(3,2)$、$B(3,-2)$，
\[
S=\frac12\cdot4\cdot3=6;
\qquad\boxed{\frac23|3\cdot3|=6}.
\]
\textbf{提醒：} $x_1=x_2$ 不代表两点重合：
它们的\textbf{纵坐标符号相反}。

\medskip
\end{examplebox}
\clearpage
\begin{examplebox}
\textbf{例5（反求另一截点的横坐标）。}

已知 $a=5,b=3$，$x_1=10$，$S_{\triangle AOB}=18$，求 $x_2$ 的可能值。
\[
18=\frac35\cdot|10x_2|=6|x_2|
\quad\Longrightarrow\quad\boxed{x_2=3\ \text{或}\ -3}.
\]
两种选择都能确定一条与两条渐近线各有一不同交点的直线。

\medskip
\textbf{例6（高价值变式：只给截线方程）。}

设 $k=b/a>0$，截线 $\ell:y=mx+n$，且
$n\ne0$、$m\ne\pm k$。
联立 $y=kx$、$y=-kx$，得到
\[
x_1=\frac{n}{k-m},\qquad x_2=-\frac{n}{k+m}.
\]
因此
\[
\boxed{S_{\triangle AOB}=\frac{k n^2}{|m^2-k^2|}}.
\]
例如 $k=\tfrac23$，$\ell:y=2x-8$，
\[
S=\frac{(2/3)\cdot64}{|4-4/9|}=\boxed{12}.
\]
\textbf{用法：}题目只给截线斜率、截距，先求 $x_1x_2$，
无需分别写出全部点坐标。$m=\pm k$ 时截线平行某条渐近线，
不满足“两条都有交点”；$n=0$ 时两个交点合为原点。

\medskip

\textbf{例7（反比例函数的跨章节认知检验）。}

初中反比例函数 $y=6/x$，取 $P=(3,2)$，设 $Q=(3,0)$。
则 $\triangle OPQ$ 面积为
\[
S_{\triangle OPQ}=\frac12\times3\times2=3=\frac{|6|}{2}.
\]
请与 C088 的 $\frac ba|x_1x_2|$ 对照：
\textbf{共同点}是面积转为横纵坐标的乘积；
\textbf{区别}是 $P$ 在反比例函数图象上，而 C088 的 $A,B$
在\emph{两条渐近线}上，不能把两个图形当作同一幅图。


\end{examplebox}
\clearpage
\begin{examplebox}
\textbf{例8（初中面积 $\to$ 高中切线：为什么出面积定值？）。}

设 $P(x_0,y_0)$ 是双曲线 $H:x^2/a^2-y^2/b^2=1$ 上一点，
其切线与两渐近线分别交于 $A,B$。\textbf{求 $S_{\triangle AOB}$。}

\begin{center}% C088 PDF teaching figure; exact coordinates, equal x/y units.
\begin{tikzpicture}[x=0.73cm,y=0.73cm,>=stealth,font=\small]
\draw[rounded corners=5pt,black!25] (-1.2,-6) rectangle (16.1,7.6);
\node[font=\bfseries] at (7.45,6.949999999999999) {切线固定横坐标乘积：$x_1x_2=a^2$，进而 $S=ab$};
\begin{scope}[shift={(0,0)}]
\node at (3.3,5.55) {切线 $x=3$：$S=ab=6$};
\draw[->,black!45] (-0.35,0)--(7.05,0) node[right] {$x$};
\draw[->,black!45] (0,-4.5)--(0,4.2) node[above] {$y$};
\draw[blue!65!black,thick] (0,0.0)--(6.65,4.433333333333334);
\draw[violet!75!black,thick,dashed] (0,-0.0)--(6.65,-4.433333333333334);
\fill[teal!18] (0,0)--(3,2.0)--(3,-2.0)--cycle;
\draw[black!65,dotted,very thick,domain=-4:4,samples=90,variable=\t] plot ({3*sqrt(1+\t*\t/4)},\t);
\draw[orange!85!black,very thick] (3,-4.55)--(3,4.55);
\fill (0,0) circle (1.7pt) node[below left=4pt] {$O$};
\fill (3,2.0) circle (1.8pt) node[above left=5pt] {$A$};
\fill (3,-2.0) circle (1.8pt) node[below left=5pt] {$B$};
\fill (3,0) circle (1.6pt) node[above left=5pt] {$P$};
\node at (3.3,-5.3) {顶点切线，$x_1x_2=9=a^2$};
\end{scope}
\begin{scope}[shift={(8.5,0)}]
\node at (3.3,5.55) {割线 $x=6$：$S=24$};
\draw[->,black!45] (-0.35,0)--(7.05,0) node[right] {$x$};
\draw[->,black!45] (0,-4.5)--(0,4.2) node[above] {$y$};
\draw[blue!65!black,thick] (0,0.0)--(6.65,4.433333333333334);
\draw[violet!75!black,thick,dashed] (0,-0.0)--(6.65,-4.433333333333334);
\fill[teal!18] (0,0)--(6,4.0)--(6,-4.0)--cycle;
\draw[black!65,dotted,very thick,domain=-4:4,samples=90,variable=\t] plot ({3*sqrt(1+\t*\t/4)},\t);
\draw[orange!85!black,very thick] (6,-4.55)--(6,4.55);
\fill (0,0) circle (1.7pt) node[below left=4pt] {$O$};
\fill (6,4.0) circle (1.8pt) node[above left=5pt] {$A$};
\fill (6,-4.0) circle (1.8pt) node[below left=5pt] {$B$};
\fill (3,0) circle (1.6pt) node[above left=5pt] {$P$};
\node at (3.3,-5.3) {任意截线未必有 $x_1x_2=a^2$};
\end{scope}

\end{tikzpicture}
\end{center}
图取 $a=3,b=2$：同尺度下，顶点切线面积 $6$，普通割线面积 $24$。

切线方程（此处引入的\textbf{新增条件}）是
\[
\frac{xx_0}{a^2}-\frac{yy_0}{b^2}=1.
\]
分别代入 $y=\pm(b/a)x$，得到
\[
x_1=\frac{a^2}{x_0-(a/b)y_0},\qquad
x_2=\frac{a^2}{x_0+(a/b)y_0}.
\]
两式相乘后，\textbf{关键是分母还原为双曲线方程}：
\[
 x_1x_2=\frac{a^4}{x_0^2-(a^2/b^2)y_0^2}
 =\frac{a^4}{a^2}=\boxed{a^2}.
\]
代入 C088，立刻有
\[
\boxed{S_{\triangle AOB}=\frac ba a^2=ab}.
\]
\textbf{考查目的：}切线条件固定了 $x_1x_2=a^2$；一般截线不保证 $S=ab$。
\end{examplebox}
